\clearpage
\section{Padding stability and copyable payloads}
\label{sec:copying-padding-predicates}
\termdef{copying-predicate-padding}
The copying machine must retain the side conditions of Definitions~10.2
and 10.3 when a finite representation is embedded at a larger allowance.
We prove this structural prerequisite for the unchanged graph definitions.
Stability is preserved and reflected; copyability is preserved. These results
neither define machine transitions nor prove execution padding invariance.
They make no change to the terms or translation of Proposition~12.2.

\subsection{Stability is independent of unused allowance}
A label is stable exactly when it denotes an instance or a passive operation.
The padding label map leaves the primitive unchanged. At allowance zero,
the proof removes the primitive wrapper; at positive allowance it preserves
the primitive branch of the sum. A box remains a box and cannot satisfy the
stable-label predicate on either side. Thus the stable-label test is equivalent
before and after one padding.
\lean{Copying/PaddingPredicates}{12}{29}{Stable.padding_label}

For a graph with one term input, the existing theorem \texttt{Stable.of\_labels}
reconstructs the entire stability predicate from its label condition for this
specific example alphabet. Its proof uses the graph's incidence coverage and
typing laws to derive both the permitted outputs and input reachability.
Consequently these conditions are not dropped in the new proof. Express
padding as the already proved uniform relabelling of the original graph;
the outer edge carrier is unchanged. Transport each edge's stable-label
condition in either direction and apply \texttt{Stable.of\_labels}. This proves
\[
  \operatorname{Stable}(\operatorname{raise}_d G)
  \quad\Longleftrightarrow\quad \operatorname{Stable}(G).
\]
\lean{Copying/PaddingPredicates}{32}{43}{Stable.raise_iff}

Induction on the number of padding steps gives the same equivalence for
iterated padding: the zero case is identity and the successor case composes
one-step equivalence with the induction hypothesis. Transport along an equality
of depth indices preserves stability by equality elimination. Combining this
transport with iterated padding proves equivalence at every specified adequate
allowance. In particular, unused allowance cannot make an unstable graph stable.
\lean{Copying/PaddingPredicates}{45}{49}{Stable.pad_iff}
\lean{Copying/PaddingPredicates}{51}{52}{Stable.castDepth_iff}
\lean{Copying/PaddingPredicates}{54}{56}{Stable.padTo_iff}

\subsection{Copyable shapes survive padding}
The lambda payload is the original serial composition of a lambda operation
and its box, before the surrounding environment distributor. Padding this
serial composition is recursively isomorphic to composing the padded
components. The primitive-padding equality identifies the lambda operation,
and the lambda-box equality identifies the original box with its body padded
once. This gives a recursive isomorphism to the lambda payload of the padded
body, with the same bound, free-variable and atom interfaces.
\lean{Copying/PaddingPredicates}{63}{67}{Copyable.raise_lambdaPayload}

Induct on the derivation of copyability. For instances, constants,
application, unary operations and binary operations, the primitive-padding
identity gives precisely the corresponding copyable constructor at the next
allowance. For a lambda payload, the padded body is arbitrary: the definition
of copyability places no stability condition on that body. Its payload is
copyable by the lambda constructor, and the inverse of the preceding
isomorphism transports copyability to the padded original payload. Finally,
if copyability was obtained by recursive renaming, pad that isomorphism and
transport the induction hypothesis along it. This proves preservation for
every constructor of the full example operation alphabet.
\lean{Copying/PaddingPredicates}{70}{79}{Copyable.raise}

Induction on the padding count gives iterated preservation. Equality
elimination handles a change of depth index, and composing these two facts
proves preservation at an arbitrary adequate allowance. No converse for
copyability is claimed here; no assumption about the behaviour of an
execution or about deterministic matching is used.
\lean{Copying/PaddingPredicates}{81}{85}{Copyable.pad}
\lean{Copying/PaddingPredicates}{87}{89}{Copyable.castDepth}
\lean{Copying/PaddingPredicates}{91}{94}{Copyable.padTo}

\subsection{Remaining operational boundary}
These theorems discharge the forward padding obligations for both side
conditions, and the reflection obligation for stability. They belong to the
structural prerequisites in Appendix~B.1. Preservation under actual machine
steps, copying-rule matching, concrete execution isomorphism and padding laws,
template robustness and the final equivalence remain to be proved. The four
outstanding work packages are therefore retained.
